\usepackage{array}
\usepackage{amssymb}
\usepackage{amsmath}
\usepackage{ntheorem}
\usepackage{epsf}
\usepackage{color}
\usepackage{xcolor}
\usepackage{graphicx}
\usepackage{epstopdf}
\usepackage{subfigure}
\usepackage{arydshln}
\usepackage{makecell}
\usepackage{mathrsfs}
\usepackage{multirow}
\usepackage{longtable}
\usepackage{threeparttable}
\usepackage{url}
\usepackage{makecell}
\usepackage{booktabs}
\usepackage{listings}
\usepackage{float}
\usepackage{diagbox}
\lstset{numbers=left, %设置行号位置
        numberstyle=\tiny, %设置行号大小
        keywordstyle=\color{blue}, %设置关键字颜色
        commentstyle=\color[cmyk]{1,0,1,0}, %设置注释颜色
        frame=single, %设置边框格式
        escapeinside=``, %逃逸字符(1左面的键)，用于显示中文
        %breaklines, %自动折行
        extendedchars=false, %解决代码跨页时，章节标题，页眉等汉字不显示的问题
        xleftmargin=2em,xrightmargin=2em, aboveskip=1em, %设置边距
        tabsize=4, %设置tab空格数
        showspaces=false %不显示空格
        title=false
}

\usepackage[bookmarks=false]{hyperref}
\hypersetup{hidelinks,
	colorlinks=true,
	allcolors=blue,
	pdfstartview=Fit,
	breaklinks=true}


\usepackage[ruled, linesnumbered]{algorithm2e}
\def\bfm#1{\mbox{\boldmath$#1$}}
\setlength{\arraycolsep}{2pt}
% \def\T{{ \mathrm{\scriptscriptstyle T} }}
\textwidth 165mm
\textheight 220mm
\numberwithin{equation}{section}

% 定义参考文献的格式
\usepackage[style=authoryear, giveninits=true, uniquename=init, backend=biber, maxbibnames=99, minbibnames=3, maxcitenames=2]{biblatex}

% ==========================================
% --- 针对期刊 (Article) 的智能定制 ---
% ==========================================
% 1. 姓名格式调整：去掉名字缩写后面的句号 (H. -> H)
\renewcommand*{\bibinitperiod}{}

% 2. 姓名格式调整：去掉姓和名之间的逗号 (Akaike, H -> Akaike H)
\renewcommand*{\revsdnamepunct}{}

% 3. 卷号调整：强制将期刊的卷号 (Volume) 设置为加粗
\DeclareFieldFormat[article]{volume}{\textbf{#1}}

% 4. 标题调整：去掉文章标题两边的默认双引号
\DeclareFieldFormat[article,inbook,incollection,inproceedings]{title}{#1}

% 5. 期刊前缀调整：去掉默认的 "In: " 字样
\renewbibmacro{in:}{%
  \ifboolexpr{test {\ifentrytype{incollection}} or test {\ifentrytype{inproceedings}}}
    {\printtext{In\space}}
    {}%
}
\DeclareFieldFormat[article]{number}{\mkbibparens{#1}}

% 3. 彻底重写“卷号+期号”的拼接逻辑：去掉它们中间的点，让它们紧紧贴在一起
\renewbibmacro*{volume+number+eid}{%
  \printfield{volume}%
  \printfield{number}% 这里直接跟着期号，不加任何点或空格
  \setunit{\addcomma\space}%
  \printfield{eid}}
\DeclareFieldFormat[article]{pages}{#1}
\DeclareNameAlias{sortname}{family-given}
\renewcommand*{\bibinitdelim}{}
\renewcommand*{\finalnamedelim}{\addspace\&\space}
\renewbibmacro*{name:andothers}{\ifboolexpr{test {\ifnumequal{\value{listcount}}{\value{liststop}}} and test \ifmorenames}{\ifnumgreater{\value{liststop}}{1}{\finalandcomma}{}\printdelim{andothersdelim}\textit{et al.}}{}}
\DeclareFieldFormat{doi}{https://doi.org/#1}
\DeclareFieldFormat{doi}{\url{https://doi.org/#1}}

\renewbibmacro*{addendum+pubstate}{%
  \printfield{addendum}%
  \iffieldundef{pubstate}
    {}
    {\setunit{\addcomma\space}% 强制将前面的标点改为逗号
     \printfield{pubstate}}%
}

% ==========================================
% --- 针对书籍 (Book) 的智能定制 ---
% ==========================================

% 1. 仅针对 book 释放卷号和版本号的默认格式
\DeclareFieldFormat[book]{volume}{#1}
\DeclareFieldFormat[book]{edition}{#1}

% 2. 带有“类型隔离”的安全书名宏
\renewbibmacro*{title}{%
  \ifboolexpr{test {\iffieldundef{title}} and test {\iffieldundef{subtitle}}}
    {}
    {\printtext[title]{%
       \printfield[titlecase]{title}%
       \setunit{\subtitlepunct}%
       \printfield[titlecase]{subtitle}}%
     % --- 核心修正开始：增加类型判断 ---
     \ifentrytype{book}
       {\setunit{\addspace}% 【如果是书籍】加空格，准备拼接括号
        \ifboolexpr{test {\iffieldundef{volume}} and test {\iffieldundef{edition}}}
          {} 
          {\printtext[parens]{% 加上小括号
             \iffieldundef{volume}{}{Vol.\addspace\printfield{volume}}%
             \ifboolexpr{not test {\iffieldundef{volume}} and not test {\iffieldundef{edition}}}
               {\addcomma\space}{}%
             \iffieldundef{edition}{}{\printfield{edition}\addspace Edition}%
           }}%
        \clearfield{volume}% 仅清除书籍的卷号，绝不碰期刊的卷号！
        \clearfield{edition}% 
       }
       {}% 【如果不是书籍】（如期刊Article），什么都不做，安全退出
     % --- 核心修正结束 ---
     \newunit}%
}

% 3. 彻底重写出版社和出版地的排列逻辑 (Publisher, Location)
\renewbibmacro*{publisher+location+date}{%
  \printlist{publisher}% 先打印出版社
  \setunit*{\addcomma\space}% 插入“逗号+空格”
  \printlist{location}% 再打印出版地
  \newunit}

% ==========================================
% --- 针对书籍章节 (Incollection) 的定制 ---
% ==========================================

% 1. 去掉页码前面的 "pp."
\DeclareFieldFormat[incollection]{pages}{#1}

% 2. 将页码强行拼接到所在书名 (booktitle) 后面，并用逗号隔开
\renewbibmacro*{booktitle}{%
  \ifboolexpr{test {\iffieldundef{booktitle}} and test {\iffieldundef{booksubtitle}}}
    {}
    {\printtext[booktitle]{%
       \printfield[titlecase]{booktitle}%
       \setunit{\subtitlepunct}%
       \printfield[titlecase]{booksubtitle}}%
     \iffieldundef{pages}
       {}
       {\setunit{\addcomma\space}% 插入逗号和空格
        \printfield{pages}% 打印页码
        \clearfield{pages}}% 关键：打印完立刻销毁，防止结尾重复打印
    }%
}

% ==========================================
% --- 针对预印本 (Preprint) 的智能定制 ---
% ==========================================
\DeclareFieldFormat[unpublished,inbook,incollection,inproceedings]{title}{#1}